\selectlanguage{ngerman}%
\frenchspacing%
\begin{ledgroupsized}[r]{120mm}%
\footnotesize%
\pstart%
\noindent\textbf{Überlieferung:}%
\pend%
\end{ledgroupsized}%
\begin{ledgroupsized}[r]{114mm}%
\footnotesize%
\pstart%
\parindent -6mm%
\makebox[6mm][l]{\textit{L}}%
Aufzeichnung:
LH XXXV 10,\;11 Bl.\;107.
Ein Blatt 8\textsuperscript{o},
Ränder zum Teil beschnitten;
ein \href{\Scan{LH_35_10_11_107wz}}{Wasserzeichen}.
Eineinhalb Seiten;
die untere Hälfte von Bl.\;\href{\Scan{LH_35_10_11_107v}}{107\,v\textsuperscript{o}} ist unbeschrieben.
Geringfügiger Textverlust am rechten Rand von Bl.\;\href{\Scan{LH_35_10_11_107r}}{107\,r\textsuperscript{o}}
\pend%
\end{ledgroupsized}%
\begin{ledgroupsized}[r]{114mm}%
\footnotesize%
\pstart%
\parindent -6mm%
\makebox[6mm][l]{\textit{Ü}}%
\textsc{G.\:W.\;Leibniz}, \textit{Deriving the laws of motion from abstract principles}. Translation and Commentary, hrsg. von \textsc{R.\;Arthur} und \textsc{O.\;Ottaviani} (\textit{Leibniz Review} XXXIV, 2024, im Erscheinen):\cite{02619}\pend%
\end{ledgroupsized}%
\vspace{5mm}%
\begin{ledgroup}%
\footnotesize%
\pstart%
\noindent%
\textbf{Datierungsgründe:}
Die Datierung der Aufzeichnung N.\;\ref{60080} beruht auf äußeren Eigenschaften des Text\-trägers %
sowie auf inhaltlichen Betrachtungen.
Siehe die ausführliche Begründung auf S.~\pageref{60080_dat}.%
\pend%
\end{ledgroup}%
\selectlanguage{latin}%
\frenchspacing%
\vspace{8mm}%
\normalsize%
\pstart%
\noindent%
\lbrack%
\href{\Scan{LH_35_10_11_107r}}%
{107\,r\textsuperscript{o}}%
\rbrack\ %
Debent motus leges%
\indexS{lex motus}
ex abstractis principiis%
\indexS{principium abstractum}
posse duci.
Et videtur maxime servire haec consideratio,%
\indexS{consideratio}
\edtext{%
quod ex motu rectilineo%
\indexS{motus rectilineus}
corporum quotcunque%
}{%
\lemma{quod}%
\Bfootnote{%
\textit{(1)}~motus semper est uniformis\indexS{motus uniformis} %
\textit{(2)}~ex motu %
\textit{(a)}~prae s %
\textit{(b)}~rectilineo corporum quotcunque%
~\textit{L}%
}}
per se discerni non potest,
in quo sit motus;
seu omnes Hypotheses succedunt.%
\indexS{hypothesis succedens}
Ergo et in consequentiis%
\indexS{consequentia}
omnes Hypotheses succedere debent.%
\indexS{hypothesis succedens}
Si nullus est modus discernendi in determinantibus,%
\indexS{modus discernendi in determinantibus}
nec debet esse in ipsis determinatis.%
\indexS{modus discernendi in determinatis}
Itaque nec in circulari
debet esse modus discernendi,%
\indexS{motus circularis}
\edlabel{LH_35_10_11_107_107r_centrip-1}%
nam firmitas corporum%
\indexS{firmitas corporis}%
\indexS{corpus firmum}
non aliunde oritur%
\indexS{origo firmitatis corporis}
quam ex vi quadam quasi centripeta,%
\indexS{vis quaedam quasi centripeta}
corpora versus se invicem agente%
\indexS{vis agens corpora versus se invicem}%
\indexS{corpora acta versus se invicem}
ob motum ambientium.%
\indexS{motus ambiens}%
\edlabel{LH_35_10_11_107_107r_centrip-2}
Et quae
\edtext{%
parum recessa%
\bracks{,}%
}{%
\lemma{parum}%
\Bfootnote{%
\textit{(1)}~recessere %
\textit{(2)}~recessa%
~\textit{L}%
}}
magna vi retrorsum aguntur,
et hoc est firmitas.%
\indexS{firmitas corporis}
Causa autem est motus%
\indexS{causa firmitatis corporis}
\edtext{%
conspiratio%
\indexS{conspiratio motus}
qui%
}{%
\lemma{conspiratio}%
\Bfootnote{%
\textit{(1)}~cum %
\textit{(2)}~quae %
\textit{(3)}~qui%
~\textit{L}%
}}
huc abiit,
nec mutari facile potest.
Nempe pro principio habendum,
quam minimum mutari.
\pend%
\pstart%
\edlabel{LH_35_10_11_107_107r_qwp-1}%
Ex principio aequivalentium hypothesium%
\indexS{principium hypothesium aequivalentium}
sequitur
eandem esse actionem et reac\-tionem,%
\edlabel{LH_35_10_11_107_107r_qwp-2}%
\indexS{actio et reactio}%
\indexS{reactio et actio}
seu ictum.
\pend%
\pstart%
Varia
\edtext{%
principia:
\newline\hspace*{7,42mm}%
(1)~%
Quam minima fit mutatio.%
\indexS{mutatio quam minima}%
}{%
\lemma{principia: (1)}%
\Bfootnote{%
\textit{(1)}~nihil fit %
\textit{(2)}~null\textlangle a\textrangle\ %
\textit{(3)}~quam minima fit mutatio.%
~\textit{L}%
}}
\newline\hspace*{7,42mm}%
(2)~%
Eadem servatur
\edtext{%
in summa%
}{%
\lemma{in summa}%
\Bfootnote{%
\textit{erg.\;L}%
}}
quantitas potentiae.%
\indexS{quantitas potentiae}%
\indexS{potentia servata}
\newline\hspace*{7,42mm}%
(3)~%
Eadem servatur quantitas directionis%
\indexS{quantitas directionis}
in
\edtext{%
summa.
\newline\hspace*{7,42mm}%
(4)~%
Motus est aliquid respectivum%
\indexS{motus respectivus}
seu omnes hypotheses%
}{%
\lemma{summa (4)}%
\Bfootnote{%
\textit{(1)}~Omnes hypotheses %
\textit{(2)}~Motus est %
\lbrack...\rbrack\ omnes hypotheses%
~\textit{L}%
}}
sunt aequivalentes.%
\indexS{hypotheses aequivalentes}%
\newline\hspace*{7,42mm}%
(5)~%
\edlabel{LH_35_10_11_107_107r_njm-1}%
Actio et reactio sunt aequales.%
\edlabel{LH_35_10_11_107_107r_njm-2}%
\indexS{actio et reactio}%
\indexS{reactio et actio}
\newline\hspace*{7,42mm}%
(6)~%
Si determinantia in se invicem abeant,%
\indexS{determinantia abeuntia in se invicem}
\ledpb%
etiam determinata in se invicem abibunt.%
\indexS{determinantia et determinata}
Hinc quies est casus motus,%
\indexS{quies ut casus motus}
aequalitas inaequalitatis.%
\indexS{aequalitas ut casus inaequalitatis}
\newline\hspace*{7,42mm}%
(7)~%
Nulla fit mutatio per saltum,%
\indexS{mutatio per saltum}
hinc omnia corpora flexilia.%
\indexS{corpora omnia flexilia}%
\indexS{corpus flexile}
\newline\hspace*{7,42mm}%
(8)~%
Communis motus%
\indexS{motus communis}
nil mutat%
\lbrack:\rbrack\
sequitur ex aequivalentia hypothesium?%
\indexS{aequivalentia hypothesium}
\pend%
\pstart%
An regula haec bona,%
\indexS{regula bona}
ut omnia serventur
quae servari possunt
\edtext{%
aliis
compa\bracks{ti}bilia
servandis.
An ergo et servanda%
}{%
\lemma{aliis}%
\Bfootnote{%
\bars{compabilia \textit{ändert Hrsg.}} servandis. %
\textit{(1)}~Quod si plura servanda \angles{mi--} %
\textit{(3)}~An ergo et servanda%
~\textit{L}%
}}
quantitas motus,%
\indexS{quantitas motus}
sed ea non potest servari,
servata quantitate potentiae;%
\indexS{quantitas potentiae}%
\indexS{potentia servanda}
et non involvit omnia.
Ope virtutis elasticae%
\indexS{virtus elastica}
revera quam
\lbrack%
\href{\Scan{LH_35_10_11_107v}}%
{107\,v\textsuperscript{o}}%
\rbrack\
minima fit mutatio.%
\indexS{mutatio quam minima}
Et hinc etiam necesse est
corpora esse flexilia.%
\pend%
\begin{edfigure}[1.75em][1.15em]%
\includegraphics[width=0.3\textwidth]%
{gesamttex/edit_VIII-4/images/LH_35_10_11_107_d_107v.pdf}%
\edcaption[0.25em]
{\hspace{0.5mm}\lbrack\textit{Fig.\:1}\rbrack}%
\end{edfigure}%
\pstart%
Ex diversis principiis
reliqua omnia semper demonstrari possunt,
sed difficile
\edtext{%
est
ultimam analysin reperire,%
\indexS{analysis ultima}
seu ostendere
quid sit omnium primum.%
\indexS{primum omnium}%
}{%
\lemma{est}%
\Bfootnote{%
\textit{(1)}~ultimum %
\textit{(2)}~ultimam analysin \lbrack...\rbrack\ omnium primum.%
~\textit{L}%
}}
Omnia tamen videntur pendere
ex hoc principio,
\edlabel{LH_35_10_11_107_107v_PID-1}%
\edtext{%
quod indiscernibilibus%
}{%
\lemma{quod}%
\Bfootnote{%
\textit{(1)}~similibus %
\textit{(2)}~indiscernibilibus%
~\textit{L}}%
}
existentibus determinantibus%
\indexS{determinantia indiscernibilia}
etiam similia sint determinata%
\indexS{determinantia et determinata}
atque discerni nequeant.%
\edlabel{LH_35_10_11_107_107v_PID-2}
Unde sequitur
quod in motu
\edtext{%
etiam post concursum%
\indexS{concursus corporum}
una Hypothesis ab alia discerni nequeat.%
\indexS{hypotheses indiscernibiles}%
}{%
\lemma{etiam}%
\Bfootnote{%
\textit{(1)}~potest %
\textit{(2)}~post concursum \lbrack...\rbrack\ discerni nequeat.%
~\textit{L}%
}}
\pend
\pstart%
Item
quod quam minima semper debet esse mutatio.%
\indexS{mutatio quam minima}
Unde non quidem sufficit
eandem quantitatem motus servari,%
\indexS{quantitas motus}
sed potentiam%
\indexS{potentia servanda}%
\lbrack,\rbrack\
\edlabel{LH_35_10_11_107_107v_ert-1}%
quia ea omnia involvit%
\lbrack,\rbrack\
nempe effectum%
\indexS{effectus potentiae}
et actionis perfectionem.%
\edlabel{LH_35_10_11_107_107v_ert-2}%
\indexS{perfectio actionis}%
\indexS{actio perfecta}
\pend%
\vspace*{6.5em}%
\pstart%
\noindent%
\centering%
*     *     *
\pend%
\newpage%
\selectlanguage{german}%
\frenchspacing%
\footnotesize%
\pstart%
\noindent%
\textls{Ausführliche Datierungsgründe:}%
\pagelabel{60080_dat}
In der titellosen und undatierten Aufzeichnung N.\;\ref{60080}, editorisch \textit{Motus leges ex abstractis principiis ductae} benannt, umreißt Leibniz eine Ableitung der Bewegungsgesetze aus allgemeineren "abstrakten Prinzipien" unter besonderer Berücksichtigung des folgenden Grundsatzes:
Sind die bestimmenden Umstände voneinander ununterscheidbar, so sind auch die bestimmten Folgen voneinander ununterscheidbar (vgl. oben, \refpassage{LH_35_10_11_107_107v_PID-1}{LH_35_10_11_107_107v_PID-2}).
Der Träger N.\;\ref{60080} weist ein \href{\Scan{LH_35_10_11_107wz}}{Wasserzeichen} auf, das noch nicht eindeutig ermittelt wurde;
somit kann nach heutigem Wissensstand nicht behauptet werden, dass es sich hierbei um die Marke einer italienischen oder österreichischen Papiermühle handele.
Eine chronologische Einordnung der Aufzeichnung ergibt sich trotzdem aus folgenden Beobachtungen (siehe hierüber \textsc{G.\:W.\;Leibniz}, \textit{Deriving the laws of motion from abstract principles}. Translation and Commentary, hrsg. von \textsc{R.\;Arthur} und \textsc{O.\;Ottaviani}, \textit{Leibniz Review} XXXIV, 2024, im Erscheinen):\cite{02619}
\pend%
\pstart%
Leibniz verwendet in der Aufzeichnung N.\;\ref{60080} eine Begrifflichkeit, die offenbar eine Lektüre von Isaac Newtons\indexN{\textls{Newton} (Neutonus), Isaac 1643\textendash1727} \textit{Philosophiae naturalis principia mathematica} (London 1687)\cite{00535} bezeugt:
etwa den Bregriff einer \textit{vis quaedam quasi centripeta} (als Ursprüng der Körperkohäsion \refpassage{LH_35_10_11_107_107r_centrip-1}{LH_35_10_11_107_107r_centrip-2}) oder die Gleichsetzung von \textit{actio} und \textit{reactio} (\refpassage{LH_35_10_11_107_107r_qwp-1}{LH_35_10_11_107_107r_qwp-2}, \refpassageohne{LH_35_10_11_107_107r_njm-1}{LH_35_10_11_107_107r_njm-2}).
Leibnizens erste bezeugte Lektüre der \textit{Principia}\cite{00535} fand – wie \textsc{Bertoloni Meli} 1993 (S.\;97)\cite{01357} nachgewiesen hat – im Herbst 1688 statt, als Leibniz sich auf dem Weg nach Italien\indexO{Italien} in Wien\indexO{Wien} aufhielt.
Aus diesem Grund dürfte N.\;\ref{60080} nicht vor dem Herbst 1688 entstanden sein.
(Siehe die Datierungsgründe des wohl ebenfalls im Herbst 1688 verfassten Entwurfs \textit{Physico-mechanicae leges ex congressu corporum ductae} \textit{LSB} VIII,\;3 N.\;68,\cite{02604} S.\;747\,f.)
Einer umfassenderen Lektüre von Newtons\indexN{\textls{Newton} (Neutonus), Isaac 1643\textendash1727} \textit{Principia}\cite{00535} widmete sich Leibniz während seines Aufenthaltes in Rom\indexO{Rom} zwischen Mai und November 1689 (vgl. \textit{Chronik}, S.\;98.)\cite{01236}
\pend%
\pstart%
Inhaltlich weist die Aufzeichnung N.\;\ref{60080} auch einen engen Zusammenhang mit Leibnizens \textit{Dynamica}\cite{01354} auf.
Dies zeigt sich zunächst wieder in der in N.\;\ref{60080} verwendeten Begrifflichkeit (vgl. etwa den Schlusssatz, \refpassage{LH_35_10_11_107_107v_ert-1}{LH_35_10_11_107_107v_ert-2}: die Potenz \textit{omnia involvit, nempe effectum et actionis perfectionem}).
Als noch bedeutender erweist sich wohl die Berührung mit der zwischen Rom\indexO{Rom} und Florenz\indexO{Florenz} abgefassten Abhandlung im Hinblick auf den besagten Grundsatz \textit{quod indiscernibilibus existentibus determinantibus etiam similia sint determinata atque discerni nequeant} (oben, \refpassage{LH_35_10_11_107_107v_PID-1}{LH_35_10_11_107_107v_PID-2}).
Nahezu identisch wird in Leibnizens \textit{Dynamica}, pars~II, sectio~III, prop.\;19 das "allgemeine Axiom" formuliert, \textit{quod quorum determinantia discerni non possunt, eorum nec discerni possint determinata} (\textit{CFP}~I, S.\;429.24--430.2; II.2, S.\;530.26--28).\cite{01354}
Demnach ist eine Entstehung der Aufzeichnung N.\;\ref{60080} vor der ersten Fassung der \textit{Dynamica}\cite{01354} unwahrscheinlich.
Diese letztere ist wohl nicht vor dem Sommer 1689 zu datieren (siehe hierüber \textsc{Robinet} 1988, S.\;81\,ff.).\cite{01240}
Folglich sollte auch N.\;\ref{60080} nicht vor dem Sommer 1689 abgefasst worden sein.
\pend%
\pstart%
Wahrscheinlich ist die Aufzeichnung N.\;\ref{60080} nicht viel später entstanden, möglicherweise auf dem Rückweg aus Italien\indexO{Italien} nach Hannover.\indexO{Hannover}
Dies könnte das im Träger anzutreffende unübliche Wasserzeichen erklären.
Ein genauer Terminus ante quem der Datierung von N.\;\ref{60080} lässt sich aber nach heutigem Wissensstand nicht ermittel.
\pend%
\selectlanguage{latin}%
\frenchspacing%
\normalsize%
